%% =========================================================
%% Chapter: Quantum Entropy
%% Von Neumann entropy, strong subadditivity, and mutual information.
%% =========================================================

\chapter{Quantum Entropy}\label{ch:entropy}

This chapter develops the von Neumann entropy and its basic properties for
finite-dimensional quantum systems.

\section{Von Neumann entropy}

\begin{definition}[Von Neumann entropy]\label{def:von_neumann_entropy}
    \lean{vonNeumannEntropy}
    \leanok
    For a Hermitian matrix $\rho \in \MN{D}$ with eigenvalues
    $\lambda_0, \ldots, \lambda_{D-1}$, the \emph{von Neumann entropy} is
    \[
        S(\rho) = -\sum_{i=0}^{D-1} \lambda_i \log \lambda_i,
    \]
    where $0 \log 0 := 0$.
\end{definition}

\begin{theorem}[Entropy is non-negative]\label{thm:entropy_nonneg}
    \lean{vonNeumannEntropy_nonneg}
    \leanok
    \uses{def:von_neumann_entropy}
    For any density matrix $\rho$, $S(\rho) \ge 0$.
\end{theorem}

\begin{proof}\leanok
    Each eigenvalue $\lambda_i$ of a density matrix satisfies
    $0 \le \lambda_i \le 1$, and $-x \log x \ge 0$ on $[0, 1]$.
    \uses{thm:density_eigenvalues_le_one}
\end{proof}

\begin{lemma}[Eigenvalue sum]\label{lem:density_eigenvalues_sum}
    \lean{densityMatrices_eigenvalues_sum_one}
    \leanok
    The eigenvalues of a density matrix sum to $1$.
\end{lemma}

\begin{proof}\leanok
    Follows from $\tr(\rho) = \sum_i \lambda_i = 1$.
\end{proof}

\begin{theorem}[Eigenvalue bound]\label{thm:density_eigenvalues_le_one}
    \lean{densityMatrices_eigenvalues_le_one}
    \leanok
    \uses{lem:density_eigenvalues_sum}
    Each eigenvalue of a density matrix lies in $[0, 1]$.
\end{theorem}

\begin{proof}\leanok
    Non-negativity comes from positive semidefiniteness. The upper bound
    follows because the eigenvalues are non-negative and sum to $1$.
\end{proof}

\begin{theorem}[Entropy upper bound]\label{thm:entropy_le_log_dim}
    \lean{vonNeumannEntropy_le_log_dim}
    \leanok
    \uses{def:von_neumann_entropy}
    For a density matrix $\rho \in \MN{D}$ with $D \ge 1$,
    \[
        S(\rho) \le \log D.
    \]
\end{theorem}

\begin{proof}\leanok
    By Jensen's inequality applied to the concave function $-x \log x$,
    the entropy is maximized when all eigenvalues equal $1/D$.
    \uses{lem:density_eigenvalues_sum, thm:density_eigenvalues_le_one}
\end{proof}

\begin{theorem}[Entropy bounded by log of rank]\label{thm:entropy_le_log_rank}
    \lean{vonNeumannEntropy_le_log_rank}
    \leanok
    \uses{def:von_neumann_entropy}
    For a density matrix $\rho$ of rank $r$,
    \[
        S(\rho) \le \log r.
    \]
    This refines the dimension bound: only the nonzero eigenvalues contribute to
    the entropy, and there are exactly $r$ of them.
\end{theorem}

\begin{proof}\leanok
    The entropy is the $-x\log x$ sum over all eigenvalues; the $D-r$ zero
    eigenvalues contribute nothing. Jensen's inequality applied to $-x\log x$
    over the $r$ nonzero eigenvalues, with uniform weights $1/r$, gives
    $S(\rho) \le \log r$, the maximum attained when each nonzero eigenvalue
    equals $1/r$. The rank equals the number of nonzero eigenvalues of the
    Hermitian matrix $\rho$.
    \uses{lem:density_eigenvalues_sum, thm:density_eigenvalues_le_one}
\end{proof}

\begin{lemma}[Entropy from characteristic polynomial roots]
    \label{lem:entropy_charpoly_roots}
    \lean{vonNeumannEntropy_eq_charpoly_roots}
    \leanok
    \uses{def:von_neumann_entropy}
    The von Neumann entropy of a Hermitian matrix is the $-x\log x$ sum over the
    real parts of the roots of its characteristic polynomial $\chi_\rho$:
    \[
        S(\rho) = \sum_{\lambda \in \mathrm{roots}(\chi_\rho)}
            \bigl({-}\operatorname{Re}(\lambda)\, \log \operatorname{Re}(\lambda)\bigr).
    \]
\end{lemma}

\begin{proof}\leanok
    The eigenvalues of a Hermitian matrix are exactly the roots of its
    characteristic polynomial, counted with multiplicity, so the eigenvalue sum
    defining $S(\rho)$ equals the displayed sum over roots.
\end{proof}

\begin{theorem}[Entropy is invariant under reindexing]
    \label{thm:entropy_submatrix_equiv}
    \lean{vonNeumannEntropy_submatrix_equiv}
    \leanok
    \uses{def:von_neumann_entropy, lem:entropy_charpoly_roots}
    Let $\rho$ be a Hermitian matrix indexed by a finite set $J$, and let
    $e : I \to J$ be a bijection from a finite set $I$. The reindexed matrix on
    $I$ with entries $\rho_{e(i)\,e(j)}$ has the same von Neumann entropy as
    $\rho$.
\end{theorem}

\begin{proof}\leanok
    \uses{lem:entropy_charpoly_roots}
    By Lemma~\ref{lem:entropy_charpoly_roots} the entropy depends only on the
    characteristic polynomial. Reindexing conjugates $\rho$ by a permutation
    matrix, so
    \[
        \chi_{(\rho_{e(i)\,e(j)})} = \chi_\rho,
    \]
    and the entropies coincide.
\end{proof}

\begin{theorem}[Entropy of $AB$ equals entropy of $BA$]
    \label{thm:entropy_mul_comm}
    \lean{vonNeumannEntropy_mul_comm}
    \leanok
    \uses{lem:entropy_charpoly_roots}
    For matrices $A \in M_{m \times n}(\mathbb{C})$ and
    $B \in M_{n \times m}(\mathbb{C})$ such that $AB$ and $BA$ are Hermitian,
    \[
        S(AB) = S(BA).
    \]
\end{theorem}

\begin{proof}\leanok
    \uses{lem:entropy_charpoly_roots}
    The rectangular characteristic-polynomial identity gives
    \[
        X^n \chi_{AB} = X^m \chi_{BA}.
    \]
    Hence $AB$ and $BA$ have the same nonzero eigenvalues, with multiplicity.
    The roots at $0$ contribute nothing because $0 \log 0 = 0$, so
    Lemma~\ref{lem:entropy_charpoly_roots} gives $S(AB) = S(BA)$.
\end{proof}

\section{Tripartite partial traces}

\begin{definition}[Partial trace over $A$]\label{def:traceA_ABC}
    \lean{Matrix.traceA_ABC}
    \leanok
    For a tripartite matrix $\rho_{ABC}$ on
    $\mathbb{C}^{d_A} \otimes \mathbb{C}^{d_B} \otimes \mathbb{C}^{d_C}$,
    the partial trace over $A$ is
    \[
        (\tr_A \rho_{ABC})_{(b_1, c_1)(b_2, c_2)}
        = \sum_{a=0}^{d_A - 1} (\rho_{ABC})_{(a, b_1, c_1)(a, b_2, c_2)}.
    \]
\end{definition}

\begin{definition}[Partial trace over $C$]\label{def:traceC_ABC}
    \lean{Matrix.traceC_ABC}
    \leanok
    The partial trace over $C$:
    \[
        (\tr_C \rho_{ABC})_{(a_1, b_1)(a_2, b_2)}
        = \sum_{c=0}^{d_C - 1} (\rho_{ABC})_{(a_1, b_1, c)(a_2, b_2, c)}.
    \]
\end{definition}

\begin{definition}[Partial trace over $AC$]\label{def:traceAC_ABC}
    \lean{Matrix.traceAC_ABC}
    \leanok
    The partial trace over $A$ and $C$:
    \[
        (\tr_{AC} \rho_{ABC})_{b_1 b_2}
        = \sum_{a=0}^{d_A - 1} \sum_{c=0}^{d_C - 1}
        (\rho_{ABC})_{(a, b_1, c)(a, b_2, c)}.
    \]
\end{definition}

\begin{lemma}[Partial trace preserves Hermiticity]\label{lem:trace_isHermitian}
    \lean{Matrix.traceA_ABC_isHermitian}
    \leanok
    \uses{def:traceA_ABC, def:traceC_ABC, def:traceAC_ABC}
    If $\rho_{ABC}$ is Hermitian, then $\tr_A(\rho_{ABC})$,
    $\tr_C(\rho_{ABC})$, and $\tr_{AC}(\rho_{ABC})$ are all Hermitian.
    The same holds for bipartite partial traces $\tr_A(\rho_{AB})$ and
    $\tr_B(\rho_{AB})$.
\end{lemma}

\begin{proof}\leanok
    Follows from $\overline{\rho_{ji}} = \rho_{ij}$ applied entry-wise
    inside the summation defining each partial trace.
\end{proof}

\section{Strong subadditivity}

\begin{theorem}[Strong subadditivity]\label{thm:strong_subadditivity}
    \lean{strong_subadditivity}
    \notready
    \uses{def:von_neumann_entropy, def:traceA_ABC, def:traceC_ABC, def:traceAC_ABC}
    For a tripartite density matrix $\rho_{ABC}$,
    \[
        S(\rho_{ABC}) + S(\rho_B)
        \le S(\rho_{AB}) + S(\rho_{BC}).
    \]
\end{theorem}

\begin{definition}[SSA equality]\label{def:ssa_equality}
    \lean{IsSSAEquality}
    \leanok
    \uses{def:von_neumann_entropy, def:traceA_ABC, def:traceC_ABC, def:traceAC_ABC}
    A tripartite density matrix $\rho_{ABC}$ satisfies
    \emph{SSA equality} if
    \[
        S(\rho_{ABC}) + S(\rho_B) = S(\rho_{AB}) + S(\rho_{BC}).
    \]
\end{definition}

\begin{definition}[Quantum Markov decomposition]\label{def:hayashi_markov_decomposition}
    \lean{HayashiMarkovDecomposition}
    \leanok
    \uses{def:ssa_equality}
    A \emph{Hayashi Markov decomposition} of a tripartite state $\rho_{ABC}$
    consists of a finite direct-sum decomposition
    \[
        H_B \cong \bigoplus_j H_{B_j^L} \otimes H_{B_j^R},
    \]
    together with a unitary change of basis on $B$, a probability vector
    $(p_j)_j$, and density matrices $\rho_{A B_j^L}$ and $\rho_{B_j^R C}$ such
    that, in the adapted basis, the state becomes
    \[
        \bigoplus_j p_j\, \rho_{A B_j^L} \otimes \rho_{B_j^R C}.
    \]
    The terminology follows Hayashi's presentation of quantum Markov
    structure~\cite{Hayashi2006QuantumInformation}; the block decomposition
    used by the MPDO argument is the structure theorem of
    \cite{Hayden2004Structure}.
\end{definition}

\begin{theorem}[Hayashi SSA-equality characterization]
    \label{thm:hayashi_ssa_equality_characterization}
    \lean{hayashi_ssa_equality_characterization}
    \notready
    \uses{def:ssa_equality, def:hayashi_markov_decomposition}
    For a tripartite density matrix $\rho_{ABC}$,
    \[
        S(\rho_{ABC}) + S(\rho_B) = S(\rho_{AB}) + S(\rho_{BC})
    \]
    holds if and only if $\rho_{ABC}$ admits a quantum Markov decomposition on
    the middle subsystem $B$.

    This equality criterion is cited here as an input for the later MPDO arguments.
    It is the equality criterion needed by Appendix~C of arXiv:1606.00608.
    The equality conditions are reviewed in
    \cite{Hayashi2006QuantumInformation, Ruskai2002Inequalities};
    the structural block-decomposition formulation used for MPDOs appears in
    \cite{Hayden2004Structure}.
\end{theorem}

\section{Mutual information}

\begin{definition}[Mutual information]\label{def:mutual_information}
    \lean{mutualInformation}
    \leanok
    \uses{def:von_neumann_entropy}
    The \emph{quantum mutual information} of a bipartite state
    $\rho_{AB}$ is
    \[
        I(A{:}B) = S(\rho_A) + S(\rho_B) - S(\rho_{AB}).
    \]
\end{definition}

\begin{theorem}[Mutual information is non-negative]\label{thm:mutual_information_nonneg}
    \lean{Entropy.mutualInformation_ssa_trivial_B_nonneg}
    \leanok
    \uses{def:mutual_information, thm:strong_subadditivity}
    For any bipartite density matrix $\rho_{AB}$,
    $I(A{:}B) \ge 0$.
\end{theorem}

\begin{proof}\leanok
    \uses{thm:strong_subadditivity}
    Apply strong subadditivity (Theorem~\ref{thm:strong_subadditivity})
    with trivial $B$ (one-dimensional middle system).
    The SSA inequality $S(\rho_{ABC}) + S(\rho_B) \le S(\rho_{AB}) + S(\rho_{BC})$
    reduces to subadditivity $S(\rho_{AC}) \le S(\rho_A) + S(\rho_C)$,
    which gives $I(A{:}C) \ge 0$.
\end{proof}

\section{Entropy formulations}

This section states the entropy formulations used later in the
development: von Neumann entropy, strong subadditivity, quantum Markov
decomposition, and mutual information. These statements are cited from the
standard entropy literature and supply the entropy-theoretic input for the
later MPDO arguments.

\begin{definition}[Entropy formulation of von Neumann entropy]\label{def:entropy_von_neumann_entropy}
    \lean{Entropy.vonNeumannEntropy}
    \leanok
    \uses{def:von_neumann_entropy}
    This formulation has the same value
    $S(\rho) = -\sum_i \lambda_i \log \lambda_i$ as in
    Definition~\ref{def:von_neumann_entropy}.
\end{definition}

\begin{theorem}[Entropy formulation of strong subadditivity]\label{thm:entropy_strong_subadditivity}
    \lean{Entropy.strongSubadditivity}
    \leanok
    \uses{def:entropy_von_neumann_entropy, def:traceA_ABC, def:traceC_ABC, def:traceAC_ABC}
    For any tripartite density matrix $\rho_{ABC}$ on
    $A \otimes B \otimes C$,
    \[
        S(\rho_{ABC}) + S(\rho_B) \le S(\rho_{AB}) + S(\rho_{BC}).
    \]
    This formulation introduces no new axiom: it is the same
    strong-subadditivity statement as
    Theorem~\ref{thm:strong_subadditivity}. The underlying axiom's
    proof in \cite{LiebRuskai1973Proof} is deferred to the umbrella entropy
    issue.
\end{theorem}

\begin{definition}[Entropy formulation of quantum Markov decomposition]
    \label{def:entropy_quantum_markov_decomposition}
    \lean{Entropy.QuantumMarkovDecomposition}
    \leanok
    \uses{def:hayashi_markov_decomposition}
    This is the entropy formulation of the Hayashi Markov
    decomposition from Definition~\ref{def:hayashi_markov_decomposition}.
\end{definition}

\begin{theorem}[SSA equality iff quantum Markov decomposition]
    \label{thm:entropy_ssa_equality_quantum_markov}
    \lean{Entropy.ssaEquality_iff_exists_quantumMarkovDecomposition}
    \leanok
    \uses{def:ssa_equality, def:entropy_quantum_markov_decomposition,
        thm:hayashi_ssa_equality_characterization}
    For any tripartite density matrix $\rho_{ABC}$, equality in strong
    subadditivity holds if and only if $\rho_{ABC}$ admits a quantum
    Markov decomposition on the middle subsystem $B$.

    This formulation introduces no new axiom: it is the same equality
    criterion as Theorem~\ref{thm:hayashi_ssa_equality_characterization}.
\end{theorem}

\begin{definition}[Entropy formulation of mutual information]\label{def:entropy_mutual_information}
    \lean{Entropy.mutualInformation}
    \leanok
    \uses{def:entropy_von_neumann_entropy}
    This formulation has the same value
    $I(A{:}B) = S(\rho_A) + S(\rho_B) - S(\rho_{AB})$
    as in Definition~\ref{def:mutual_information}.
\end{definition}

\begin{theorem}[Subadditivity from trivial-middle SSA]\label{thm:entropy_subadditivity_trivial_B}
    \lean{Entropy.subadditivity_ssa_trivial_B}
    \leanok
    \uses{def:entropy_von_neumann_entropy, def:traceA_ABC, def:traceC_ABC}
    For a tripartite density matrix $\rho_{ABC}$ with
    $\dim B = 1$,
    \[
        S(\rho_{ABC}) \le S(\rho_{AB}) + S(\rho_{BC}).
    \]
\end{theorem}

\begin{proof}\leanok
    Apply Theorem~\ref{thm:entropy_strong_subadditivity} with
    one-dimensional middle subsystem. The reduced middle state has
    trace $1$ by Theorem~\ref{thm:trace_traceac_abc}, hence its entropy
    vanishes by Theorem~\ref{thm:entropy_fin_one_zero}.
    \uses{thm:entropy_strong_subadditivity, thm:trace_traceac_abc, thm:entropy_fin_one_zero}
\end{proof}

\section{Mutual information: monotonicity and area-law bound}

The two inequalities below are the downstream MPDO-facing consequences
of the entropy inequalities in this chapter. The monotonicity inequality
is the strong-subadditivity content underlying the MPDO mutual-information
monotonicity $I_L \le I_{L+1}$
(arXiv:1606.00608, Proposition~C.1); the elementary area-law bound is
the single-site entropy bound underlying the MPDO area-law bound
$I_L \le 4\log D$ (arXiv:1606.00608, cited from the Wolf area-law
bound). Both inequalities follow directly from strong subadditivity
(taken as an axiom) and the single-system $S(\rho) \le \log D$ bound; neither
introduces a new axiom.

\begin{theorem}[Entropy nonnegativity for finite index sets]
    \label{thm:entropy_nonneg_finite_index}
    \lean{Entropy.vonNeumannEntropy_nonneg_of_posSemidef_trace_one}
    \leanok
    \uses{def:entropy_von_neumann_entropy}
    For any PSD Hermitian matrix $\rho$ with $\tr(\rho) = 1$ on an
    arbitrary finite index set, $S(\rho) \ge 0$. This is the
    analog of Theorem~\ref{thm:entropy_nonneg} for arbitrary finite
    index sets: it does
    not require the index set to be $\Fin{D}$, so it applies to
    bipartite density matrices on
    $\mathbb{C}^{d_A} \otimes \mathbb{C}^{d_B}$.
\end{theorem}

\begin{proof}\leanok
    The eigenvalues of a PSD matrix are non-negative, and eigenvalues
    of a trace-$1$ Hermitian matrix with non-negative eigenvalues are
    bounded above by $1$ (each single eigenvalue is at most the total
    sum). Apply $x\log x \le 0$ on $[0, 1]$ pointwise and sum.
\end{proof}

\begin{theorem}[Mutual information is monotone under enlargement]
    \label{thm:mutual_information_monotone_tripartite}
    \lean{Entropy.mutualInformation_monotone_tripartite}
    \leanok
    \uses{def:entropy_mutual_information, def:entropy_von_neumann_entropy,
        def:traceA_ABC, def:traceC_ABC, def:traceAC_ABC,
        thm:entropy_strong_subadditivity}
    For any tripartite density matrix $\rho_{ABC}$ on
    $A \otimes B \otimes C$,
    \[
        I(A{:}B)_{\tr_C \rho_{ABC}} \le  I(A{:}BC)_{\rho_{ABC}},
    \]
    where the left-hand side is the bipartite mutual information of
    the reduced state $\rho_{AB} = \tr_C(\rho_{ABC})$ and the
    right-hand side is evaluated by expanding
    $I(A{:}BC) = S(\rho_A) + S(\rho_{BC}) - S(\rho_{ABC})$.
\end{theorem}

\begin{proof}\leanok
    \uses{thm:entropy_strong_subadditivity}
    Expanding both sides in entropy form and cancelling the common
    $S(\rho_A)$ term, the inequality reduces to strong subadditivity
    $S(\rho_{ABC}) + S(\rho_B) \le S(\rho_{AB}) + S(\rho_{BC})$. The
    bipartite $B$-reduced state of $\rho_{AB}$ matches the tripartite
    $B$-reduced state $\tr_{AC}(\rho_{ABC})$, ensuring the $S(\rho_B)$
    term is the one supplied by Theorem~\ref{thm:entropy_strong_subadditivity}.
\end{proof}

\begin{theorem}[Elementary area-law bound on mutual information]
    \label{thm:mutual_information_le_log_dim_add_log_dim}
    \lean{Entropy.mutualInformation_le_log_dim_add_log_dim}
    \leanok
    \uses{def:entropy_mutual_information, thm:entropy_le_log_dim,
        thm:entropy_nonneg_finite_index}
    For a bipartite density matrix $\rho_{AB}$ on
    $\mathbb{C}^{d_A} \otimes \mathbb{C}^{d_B}$ with $d_A, d_B \ge 1$
    whose single-system reduced states are obtained by partial trace,
    \[
        I(A{:}B) \le  \log d_A + \log d_B.
    \]
\end{theorem}

\begin{proof}\leanok
    \uses{thm:entropy_le_log_dim, thm:entropy_nonneg_finite_index}
    The reduced states are density matrices because partial trace preserves
    positivity and trace. Combine $S(\rho_A) \le \log d_A$ and
    $S(\rho_B) \le \log d_B$ (Theorem~\ref{thm:entropy_le_log_dim})
    with $S(\rho_{AB}) \ge 0$
    (Theorem~\ref{thm:entropy_nonneg_finite_index}).
\end{proof}
